\subsection{3. Schaltbarkeit und Arbeit werden
mitgerechnet}\label{schaltbarkeit-und-arbeit-werden-mitgerechnet}

Es gilt

\[\adjustbox{max width=\linewidth}{$\displaystyle 
Z_bH_3Z_b=\Delta B-gX_3,\qquad[Z_b,H_3]\ne0\quad(g\ne0).
$}\]

Jedes Wort ausschließlich aus dem fixen H bzw. seinen freien
Entwicklungen kommutiert dagegen mit H. Deshalb kann Zb nicht als
zustandsunabhängiges Wartewort des unverändert laufenden H ausgegeben
werden.

Eine konkrete \textbf{bedingte} finite Umsetzung im unabhängig
schaltbaren Tier \{X,Nb\} ist: Beide Arme entwickeln zunächst τ mit H;
für dieselbe Dauer T\_p=π/Δ wird die Paarwechselwirkung pausiert. Im
Referenzarm ist zusätzlich der Nb-Term pausiert, also H\_ref=0. Im
Impulsarm läuft H\_ctrl=ΔNb. Danach folgt in beiden Armen erneut τ mit
H. Die Gesamtdauer ist identisch, und die mittleren Entwicklungen sind
exakt I bzw. Zb. Das verlangt entsprechende Nullstellungen oder
Kompensationen beider Kontrollkoeffizienten. Eine Quelle, die lediglich
das feste H liefert, hat diese Möglichkeit noch nicht geliefert. Eine
idealisierte Instantanphase ohne diese Erläuterung würde die
Kontrollressource verstecken.

Auch die Energieänderung ist nicht null. Der Ausgangszustand und seine
freie Vorentwicklung haben Erwartungsenergie null. Der Impuls deponiert

\[\adjustbox{max width=\linewidth}{$\displaystyle 
\Delta E_{\rm Phase}=\frac\Delta4(1-d^2)
=\frac\Delta{54}\quad\text{bei }g/\Delta=1/20.
$}\]

Die Kontrolle muss diese Arbeit bereitstellen. Ladungserhaltung ersetzt
keine Energiebilanz des Steuerapparats.

\subsection{4. Ein echtes, die Paarkohärenz erhaltendes
Aufzeichnungsinstrument}\label{ein-echtes-die-paarkohuxe4renz-erhaltendes-aufzeichnungsinstrument}

Man gewähre einen zweistufigen, ladungs- und G-trivialen Zeiger im
Zustand \textbar0⟩R sowie einen schaltbaren gemeinsamen Koppler. Die
Isometrie

\[\adjustbox{max width=\linewidth}{$\displaystyle 
\boxed{V_{\rm rec}\psi=(I-B)\psi\otimes|0\rangle_R
+B\psi\otimes|1\rangle_R}
$}\]

bleibt in \(\mathcal K_3\otimes\mathbb C^2\), Dimension sechs. Sie ist
die Einschränkung des vollständigen unitären Kopplers

\[\adjustbox{max width=\linewidth}{$\displaystyle 
C_R=(I-B)\otimes I+B\otimes X_R
=\exp\!\left[-\frac{i\pi}{2}B\otimes(I-X_R)\right].
$}\]

Bei fortgesetzter Definition mit dem ganzzahligen Nb auf dem gesamten
Fockraum zeichnet dieser Koppler Bosonparität auf; im vorliegenden
N=2-Raum ist diese identisch mit Nb=0 oder 1. Die gemeinsame Kopplung
erhält N und die innere Gruppe, ist aber als Zeigerkopplung eine
\textbf{zusätzliche Ressource}. Aus der Verfügbarkeit eines
CP-Phasenkanals folgt nicht automatisch dessen kohärent kontrollierte
Umsetzung. Bei einer endlichen Messzeit müssen die freie Entwicklung und
der Vergleichsarm ebenso kontrolliert werden wie beim Phasenpuls; hier
wird die angegebene Isometrie als Instrument gewährt.

Alle Fermionenpaare erhalten dasselbe Zeigerlabel 0. Ihre Amplituden und
Interferenzen werden nicht einzeln aufgezeichnet. Es gilt Vrec†Vrec=I3:
Die Abbildung erhält sämtliche gemeinsamen Skalarprodukte. Das bedeutet
keine Gleichheit der zukünftigen reduzierten Systemstatistik mit dem
Prozess ohne Recorder. Der Recorder und sein Zeiger sind ein anderer,
vollständig mitgerechneter Prozess.

Wenn der Zeiger ignoriert wird, lautet der \textbf{wirkliche
Zustandsupdate}

\[\adjustbox{max width=\linewidth}{$\displaystyle 
\mathcal D_B(\rho)=(I-B)\rho(I-B)+B\rho B
=\frac12(\rho+Z_b\rho Z_b).
$}\]

Der Gesamtprozess ist CPTP. Er hat auf den Empfänger unmittelbar keinen
Effekt, da beide Krausoperatoren mit n5 kommutieren. Die
Paar-Boson-Kohärenz des isolierten Systems geht dabei verloren; die
Kohärenz im Paarblock bleibt erhalten. Dies ist eine gezielt gebaute
grobe Lüders-Messung, kein feines Paarlabelinstrument mit anschließend
umbenannten Ausgängen.

Wird der Recorder am selben mittleren Zeitpunkt eingesetzt und der
Zeiger bei der abschließenden n5-Statistik nicht ausgewählt, ergibt sich
exakt

\[\adjustbox{max width=\linewidth}{$\displaystyle 
p_{\rm rec}=\frac12(p_{\rm frei}+p_{Z_b}),\qquad
\boxed{\delta_{\rm rec}=\frac12\delta_{b\to5}}.
$}\]

Am Prüfpunkt beträgt die Differenz +0.0001552085135631517. Jeder
Zeigerzweig und sein Gewicht wird mitgeführt; es gibt keine
Postselektion. Die Aufzeichnung deponiert im System die Energie
Δ(1−d²)/8=Δ/108. Ein anfänglich bereiter, ladungsneutraler Zeiger ist
daher keine kostenlose Energiequelle.

Allgemeiner multiplizieren Zeigerzustände mit reellem Überlapp η
zwischen 0 und 1 die reduzierte Paar-Boson-Kohärenz mit η. Die
Empfängerwirkung wird \((1-\eta)\delta_{b\to5}/2\). Vollständig gleiche
Zeigerzustände (η=1) bewahren diese Kohärenz, enthalten aber keine
Information über die Zusammensetzung. Perfekt unterscheidbare Zustände
(η=0) liefern den gerade berechneten Recorder. Die
Informations-/Rückwirkungsgrenze bleibt sichtbar.

Die terminale n5-Messung besitzt weiterhin die
v1.6.8-Instrumentengrenze: Ihr Effekt ist exakt auf K3 komprimierbar,
ihre Zustandsrückwirkung erhält K3 nicht. Eine Wiederverwendung nach
dieser Endmessung wird nicht behauptet. Der hier neue Nb-Recorder selbst
bleibt dagegen einschließlich Zeiger im angegebenen sechs-dimensionalen
Raum.

\subsection{5. Warum reine innere Symmetrieimpulse diese Rolle nicht
übernehmen}\label{warum-reine-innere-symmetrieimpulse-diese-rolle-nicht-uxfcbernehmen}

Sei V ein tatsächlicher innerer Gruppenimpuls. Dann {[}V,H{]}=0. Wenn
sein Senderbereich anfangs mit einem Empfängeroperator O kommutiert,
gilt

\[\adjustbox{max width=\linewidth}{$\displaystyle 
[V,O]=0\quad\Longrightarrow\quad
[V,O(t)]=[V,e^{itH}Oe^{-itH}]=0
$}\]

für alle Zeiten. Folglich ist die unbedingte Empfängeränderung exakt
null, für jeden Anfangszustand. Derselbe Schluss gilt für einen
spurtreuen Krausprozess, dessen sämtliche Krausoperatoren diese beiden
Kommutationsbedingungen erfüllen. Das ist kein Schluss nur aus einer
verschwindenden linearen Antwort auf einem einzelnen Zustand.

Die Voraussetzung ist Operator-Kommutation. Wenn V den Empfänger bereits
direkt umzeichnet, oder nur zufällig sein Anfangserwartungswert gleich
bleibt, folgt dieser Satz nicht. Ein solcher direkter Zugriff wäre aber
kein verzögertes Signal zwischen zuerst kommutierenden
Operationsbereichen.

Zb ist \textbf{G-invariant}, aber kein solcher \textbf{G-Impuls}: Es
kommutiert mit der Gruppenwirkung, jedoch nicht mit H. Es greift in die
zwei Vorkommen desselben inneren Typs ein, nämlich Paar und Boson. Genau
dieser Unterschied zwischen einer inneren Gruppenaktion und einer
invarianten Operation auf den Multiplizitäten trägt den neuen Versuch.

\subsection{6. Verifikation und verbleibender
Herkunftsschritt}\label{verifikation-und-verbleibender-herkunftsschritt}

Für den parallelen Quellvergleich ist zusätzlich ein
dimensionsunabhängiger kleiner Blocksatz geprüft: Bei
\(H_C=\left(\begin{smallmatrix}0&gC^\dagger\\gC&\Delta I\end{smallmatrix}\right)\)
und einem normierten Eingang v auf der Bosonseite ist die erste und
dritte Ableitung der Boson-Überlebenswahrscheinlichkeit null, während
\(p''_{N_b}(0)=-2g^2\langle v,CC^\dagger v\rangle\). Somit gilt
\(p_{N_b}(t)=1-g^2\langle CC^\dagger\rangle t^2+O(t^4)\). Dieser reine
Blocksatz verbindet dieselbe Nb-Auslese mit einem
Kopplungs-Gramoperator; er leitet die Eingangspräparation oder die
präzise Schätzung der Anfangskrümmung nicht her.

\codeword{verify\_protocol.py} prüft Eingabepins, W und den identischen
nativen Adapter, die fünfdimensionale Kontrollalgebra, beide
Phasenwahrscheinlichkeiten, die Vorzeichenregion, Energieänderungen, das
echte Zeigerinstrument und die allgemeine Kommutatorinduktion. Ein
unabhängiger Matrixexponential-Replay bestätigt den dreidimensionalen
Phasenversuch und den sechs-dimensionalen Aufzeichnungsprozess. Normal
und \codeword{-OO} werden mit expliziten Guards und Warnungen als Fehler
ausgeführt. Zahlen und vollständige Prüfbedingungen stehen in
\codeword{protocol\_normal.json} und
\codeword{protocol\_optimized.json}. Die abschließenden Läufe enthalten
\textbf{92 exakte und vier numerische Prüfbedingungen}. Sie sind
byte-identisch; der optimierte Lauf wurde zusätzlich aus einem fremden
Arbeitsverzeichnis gestartet. Beide Ergebnisdateien besitzen SHA-256
\codeword{5f6dbae2f5a62204cf5a03971389e301268678f3b52268e40c609d32b93141ae}.

\textbf{Erreicht:} Ein zuvor separat nötiger Fermionmodenimpuls wird
durch den bereits dokumentierten Kontrolltyp Nb ersetzt. Ein kleiner
gemeinsamer Recorder bewahrt die relevante Paarkohärenz und liefert
einen eigenen unbedingten kausalen Effekt mit vollständiger Ladungs- und
Arbeitsangabe.

\textbf{Weiter offen:} Compilerseitige Schaltbarkeit von X/Nb,
ursprüngliche Präparation von p0, kalibrierter n5-Detektor,
bereitgestellter Zeiger, sein kohärenter Koppler und die Quelle der
Steuerarbeit. Das konkrete Ergebnis reduziert die Zahl unabhängiger
Operatortypen; es macht daraus noch keinen autonomen Universumsprozess
und keine räumliche Welt.

\begin{center}\rule{0.5\linewidth}{0.5pt}\end{center}

\section{Teilbericht C: minimaler Lorentzadapter und
Kinetiktest}\label{teilbericht-c-minimaler-lorentzadapter-und-kinetiktest}

\section{v1.6.9: Ein minimaler Lorentz-Auslesekanal mit
Impuls}\label{v1.6.9-ein-minimaler-lorentz-auslesekanal-mit-impuls}

\subsection{Ergebnis}\label{ergebnis}

Die gemischthändige Feldfrage wurde konstruktiv weitergelöst. Für den
sechsdimensionalen Träger \codeword{(1,1/2)} existiert mit \textbf{genau
einem zusätzlichen Lorentzvektor p} eine explizite Auslese in einen
linken Weylspinor. Innerhalb dieser Polynomialklasse ist sie bis auf
Normierung eindeutig; ohne p ist sie unmöglich. Ihr Rang ist bei jedem
von null verschiedenen p exakt zwei, auch auf dem Lichtkegel.

Die zusätzliche Struktur ist dabei wesentlich. Für zeitartiges p lässt
sich der frühere gleichhändige Projektor kovariant auf den
gemischthändigen Träger übertragen. Für lichtartiges p bleibt die
Auslese wohldefiniert, aber eine nur von p abhängige kovariante rechte
Inverse existiert nicht. Ein zweiter zeitartiger Bezugsvektor q
ermöglicht eine solche Inverse.

Der vorgezogene Kinetiktest liefert einen weiteren konkreten Befund:
Wird das Komposit aus einem freien masselosen selbstdualen Feld und
einem freien masselosen Weylfeld gebildet, verschwindet diese
Ableitungsauslese auf deren Bewegungsgleichungen identisch. Die richtige
Darstellung allein erzeugt somit noch keinen nichtverschwindenden
Weyl-Pol.

\textbf{Status:} Konstruktion und Gegenprüfungen exakt; native Herkunft
von p, q, Feldadapter und Kinetik offen. Es wurde kein räumlicher Träger
als aus W hergeleitet ausgegeben. Alle neuen Dateien liegen
ausschließlich in diesem Arbeitsstrang.

\subsection{1. Präzise zusätzliche
Eingabe}\label{pruxe4zise-zusuxe4tzliche-eingabe}

Seien V und V̄ die beiden fundamentalen Weylräume. Der neue Vertrag setzt
explizit

\[\adjustbox{max width=\linewidth}{$\displaystyle 
 T\in\operatorname{Sym}^2V\otimes\bar V=(1,1/2),
 \qquad P\in V\otimes\bar V=(1/2,1/2).
$}\]

Für reellen Impuls verwenden wir

\[\adjustbox{max width=\linewidth}{$\displaystyle 
 P=\begin{pmatrix}p_0+p_3&p_1-ip_2\\p_1+ip_2&p_0-p_3\end{pmatrix},
 \qquad \det P=p_0^2-\mathbf p^2.
$}\]

Unter M∈SL(2,C) gilt \codeword{P→M P M†} und \codeword{T→(Sym²M⊗M*)T}.
Die letztere Gleichung ist ein \textbf{angenommener Lorentzadapter} für
das Komposit. Der interne Spin(10)-Spinor und der Lorentzspinor sind
verschiedene Faktoren.

Eine echte Feldadjunktion vertauscht linke und rechte Weyldarstellungen.
Ein bloßer Fock-Erzeuger ist dagegen noch kein Feld mit diesen Indizes:
Eine Weylfeldentwicklung enthält sowohl Vernichtungs- als auch
Erzeugungsterme mit passenden Spinorwellenfunktionen. Beide Unterschiede
bleiben erhalten; siehe
\href{https://arxiv.org/pdf/0812.1594\#page=24}{Dreiner, Haber und
Martin, Abschnitte 2 und 3.1}.

\subsection{2. Explizite Auslese und
Minimalität}\label{explizite-auslese-und-minimalituxe4t}

Mit ε als alternierender Form lautet die Auslese

\[\adjustbox{max width=\linewidth}{$\displaystyle 
 (C_PT)^\alpha=\sum_{\beta,\dot\gamma,\delta,\dot\epsilon}
 T^{\alpha\beta\dot\gamma}\,
 \varepsilon_{\beta\delta}\varepsilon_{\dot\gamma\dot\epsilon}
 P^{\delta\dot\epsilon}.
$}\]

In der bisherigen unnormierten Reihenfolge
\codeword{(00,dot0),(00,dot1),(01,dot0),(01,dot1),(11,dot0),(11,dot1)}
und für \codeword{P=[[a,b],[c,d]]} wird daraus die unmittelbar
ausführbare Matrix

\[\adjustbox{max width=\linewidth}{$\displaystyle 
 C_P=\begin{pmatrix}
 d&-c&-b&a&0&0\\0&0&d&-c&-b&a
 \end{pmatrix}.
$}\]

Der neue Prüfer zeigt sämtliche sechs infinitesimalen Lorentz-
Intertwining-Gleichungen und vier endliche Transformationen
einschließlich Boost und komplexer Nullrotation. Insbesondere gilt exakt

\[\adjustbox{max width=\linewidth}{$\displaystyle 
 C_{MPM^\dagger}(\operatorname{Sym}^2M\otimes\bar M)=M C_P.
$}\]

Ohne Impuls hat das lineare Intertwining-System auf den zwölf möglichen
Auslesekoeffizienten Rang zwölf: Der Kern ist null. Mit einem
Impulsfaktor hat es auf 48 Koeffizienten Rang 47: Der Lösungsraum ist
eindimensional. Dies bestätigt die Darstellungssumme

\[\adjustbox{max width=\linewidth}{$\displaystyle 
 (1/2,1/2)\otimes(1,1/2)
 =(3/2,1)\oplus(3/2,0)\oplus(1/2,1)\oplus(1/2,0).
$}\]

Die Zielkomponente tritt genau einmal auf. Die Minimalität gilt für
\textbf{eine lokale lineare Auslese, polynomial in einem einzigen
unabhängigen Impuls bis zum angegebenen Grad}. Zwei unabhängig
zugängliche Konstituentenimpulse oder zusätzliche Felder erweitern den
Vertrag.

\subsection{3. Exakte Ränge und Kerne}\label{exakte-ruxe4nge-und-kerne}

Vier Zweierminoren von C\_P sind \codeword{d²,c²,b²,a²}. Deshalb ist der
Rang für jede von null verschiedene komplexe Matrix P exakt zwei,
unabhängig von \codeword{det P}. Alle folgenden Aussagen sind
symbolische Identitäten, keine numerischen Rangschätzungen.

{\def\LTcaptype{none} % do not increment counter
\begin{longtable}[]{@{}
  >{\raggedright\arraybackslash}p{(\linewidth - 6\tabcolsep) * \real{0.2143}}
  >{\raggedleft\arraybackslash}p{(\linewidth - 6\tabcolsep) * \real{0.2857}}
  >{\raggedleft\arraybackslash}p{(\linewidth - 6\tabcolsep) * \real{0.2857}}
  >{\raggedright\arraybackslash}p{(\linewidth - 6\tabcolsep) * \real{0.2143}}@{}}
\toprule\noalign{}
\begin{minipage}[b]{\linewidth}\raggedright
Impuls
\end{minipage} & \begin{minipage}[b]{\linewidth}\raggedleft
Rang C\_P
\end{minipage} & \begin{minipage}[b]{\linewidth}\raggedleft
Kerndimension
\end{minipage} & \begin{minipage}[b]{\linewidth}\raggedright
Kovarianter Rückweg nur mit P
\end{minipage} \\
\midrule\noalign{}
\endhead
\bottomrule\noalign{}
\endlastfoot
P=0 & 0 & 6 & Nein \\
Zeitartig, P²\textgreater0 & 2 & 4 & Explizit vorhanden \\
Lichtartig, P≠0 und P²=0 & 2 & 4 & Exakt ausgeschlossen \\
Raumartig, P²\textless0 & 2 & 4 & Algebraisch vorhanden; keine positive
massive Interpretation behauptet \\
\end{longtable}
}

Bei \codeword{P=m I} sind die Kerngleichungen \codeword{T0+T3=0},
\codeword{T2+T5=0}. Eine Basis ist \codeword{e1,e4,e0−e3,e2−e5}. Für den
Lichtkegelrepräsentanten \codeword{P=diag(1,0)} lautet die Auslese
\codeword{(T3,T5)} und der Kern \codeword{span(e0,e1,e2,e4)}. Die
Indizes sind hier nullbasiert.

\subsubsection{Zeitartiger Rückweg}\label{zeitartiger-ruxfcckweg}

Eine kovariante Einsetzung ist

\[\adjustbox{max width=\linewidth}{$\displaystyle 
 (J_P\chi)^{\alpha\beta\dot\gamma}
 =\chi^\alpha P^{\beta\dot\gamma}+\chi^\beta P^{\alpha\dot\gamma},
 \qquad
 J_P=\begin{pmatrix}
 2a&0\\2b&0\\c&a\\d&b\\0&2c\\0&2d
 \end{pmatrix}.
$}\]

Die tragende Identität lautet

\[\adjustbox{max width=\linewidth}{$\displaystyle 
 C_PJ_P=3\det(P)I_2.
$}\]

Für \codeword{det P≠0} ergeben sich damit

\[\adjustbox{max width=\linewidth}{$\displaystyle 
 R_P=\frac{J_P}{3\det P},\qquad
 \Pi_P=\frac{J_PC_P}{3\det P},\qquad
 C_PR_P=I_2,\quad\Pi_P^2=\Pi_P,\quad\operatorname{rank}\Pi_P=2.
$}\]

Dies verbindet die neue Rechnung unmittelbar mit dem eingefrorenen
gleichhändigen Adapter C\_alt,R\_alt aus v1.6.8. Für
\codeword{S\_P=I\_3⊗(P εᵀ)} prüft das Programm exakt

\[\adjustbox{max width=\linewidth}{$\displaystyle 
 C_P=C_{\rm alt}S_P,\quad
 R_P=S_P^{-1}R_{\rm alt},\quad
 \Pi_P=S_P^{-1}\Pi_{\rm alt}S_P.
$}\]

Der Impuls liefert also die fehlende Umwandlung des gepunkteten Faktors.
Auf dem zukünftigen zeitartigen Kegel existiert zusätzlich ein
positives, impulsabhängiges Tensor-Skalarprodukt, bezüglich dessen Π\_P
selbstadjungiert ist. Es entsteht aus den positiven Faktoren
\codeword{P⁻¹,P⁻¹,P̄⁻¹}; die Kovarianz und Selbstadjungiertheit sind
geprüft. Dies ist ein positives Skalarprodukt auf den vorgeschriebenen
Impulsfasern, noch keine positive Feldkinetik.

